\begin{enumerate}
    \item Compléter le tableau ci-dessous en indiquant les formules des composés
          ioniques solides obtenus en associant le cation d'une ligne et l'anion
          d'une colonne.

          \begin{minipage}{\linewidth}
              \begin{table}[H]
              \centering
              %\setlength{\arrayrulewidth}{0.8pt}
              %\setlength{\tabcolsep}{18pt}
              \renewcommand{\arraystretch}{1.5}
              \begin{tabularx}{\textwidth}{|p{3.5cm}|C|C|C|}
                  \cline{2-4}
                  \multicolumn{1}{c|}{} & chlorure \ce{Cl-} &
                  nitrate \ce{NO$_{3}^{-}$} & sulfate \ce{SO$_{4}^{2-}$} \\
                  \hline
                  sodium \ce{Na+} &  &  &  \\
                  \hline
                  fer(III) \ce{Fe^{3+}} &  &  &  \\
                  \hline
                  cuivre(II) \ce{Cu^{2+}} &  &  &  \\
                  \hline
                  ammonium \ce{NH$_{4}^{+}$} &  &  &  \\
                  \hline
              \end{tabularx}
          \end{table}
          \end{minipage}
          
    \item Écrire les équations des réactions de dissolution de ces composés dans
          l'eau.
    \item Quelle masse de chlorure de sodium doit-on dissoudre pour obtenir
          $\qty{100}{\mL}$ de solution de concentration égale à
          $\qty{0.020}{\mol\per\L}$~?
    
          Que valent les concentrations $[\ce{Cl^{-}(aq)}]$ et
          $[\ce{Na^{+}(aq)}]$ dans cette solution~?
\end{enumerate}

